% SPDX-FileCopyrightText: Copyright (c) 2023 Yegor Bugayenko
% SPDX-License-Identifier: MIT

\documentclass{article}
\usepackage[nominutes,template,scheme=dark]{ppt-slides}
\usepackage[static]{clicks}
\usepackage{multicol}
\usepackage[synonym=hs,colored,color=orange]{handstroke}
\usepackage{../../lecture-notes/notes}
\usepackage{svg}
\usepackage{tikz}
  \usetikzlibrary{shadows}

\graphicspath{{../../faces/}}

\newcommand\citch[1]{
  \plush{
    \begin{textblock}{10}[0.5,0.5](8,8)%
      \color{.}
      \centering #1%
    \end{textblock}}}

\newcommand\pitch[6]{
  \lnPitch{
    \pptBanner{#1}%
    \vspace*{2em}%
    {\large\enquote{#2}}\par%
    \href{#3}{#4}}}

\RemoveFromHook{enddocument}[notes]

\begin{document}

\newcommand*{\thedate}{October 17th, 2026}
\newcommand*{\thetitle}{Open Source}
\newcommand*{\thesubtitle}{Our Typical Mistakes When Using AI}
\newcommand*{\theauthor}{Yegor Bugayenko}

\pptLeft{\thedate}
\pptRight{@yegor256}

\citch{
  \pptTitle{\thetitle}{\thesubtitle}
  \par
  \theauthor
  \\
  \thedate
  }

\lnPitch{
  \color{.}
  According to GitHub leadership, the platform processed roughly 1 billion commits in 2025. In 2026, that number is expected to exceed 14 billion commits.
  \par
  \small
  --- @naro7217 \href{how-github-handles-billions-of-commits-while-ai-coding-tools-flood-the-platform-aa02577be670}{on Medium}, Jul 24, 2026}

\citch{%
  \qrcode[height=12em]{https://t.me/yegor256news}
  \par
  \href{https://t.me/yegor256news}{@yegor256news}
  \par
  }

\end{document}
